\paragraph{Graduation, unité et automorphismes.}
\label{inc-ampl-automorfismos}
La construction de Frenkel--Lepowsky--Meurman s'applique au réseau
$N_\alpha$ une fois ses propriétés vérifiées. Elle incorpore l'espace
des oscillateurs, l'extension centrale du réseau, le produit de vertex
et le secteur tordu par le relèvement de la négation. Le résultat est
l'algèbre $V^\natural_{\rm HMT}$ de \eqref{exc:flm}, dont le groupe
d'automorphismes est isomorphe au Monstre. Son domaine d'action est ainsi
fixé par une structure algébrique munie d'une graduation et d'un produit. Le théorème
de Moonshine de Borcherds détermine ensuite les traces graduées
avec insertion de ces automorphismes \cite{flm,borcherds}.

L'unité acquiert à ce niveau une réalisation précise : la composante
de poids zéro est la droite du vide et fournit le coefficient $1$ de
$q^{-1}$ dans le caractère gradué. Au niveau cylindrique, les
applications de raffinement conservent la fonction constante par
précomposition. Chaque assertion de conservation possède donc son
objet et son application. Le développement des preuves rend visible cette
continuité de construction, depuis les opérations arithmétiques et la
mémoire jusqu'à l'incidence, à la métrique et aux produits de vertex.

Cette articulation situe $K$ et la clôture de $\alpha$ dans l’argument
principal. $K$ conserve réversiblement une coordonnée entière de
l’histoire ; la normalisation de $\alpha$ compose les registres régionaux
avec cette coordonnée ; la configuration signée détermine un drapeau ; et
le drapeau marqué sélectionne la modification réticulaire sur laquelle
Moonshine est réalisé. La proposition~\ref{exc:selector-bandera},
le lemme~\ref{exc:origen} et les théorèmes~\ref{exc:teorema-leech}
et~\ref{exc:moonshine} établissent ces étapes. Les équivalences
réversibles, les quotients d’information et les réalisations qui ajoutent
une métrique ou une structure algébrique conservent ainsi leurs domaines respectifs.
La relation entre structure, équation et valeur s’exprime par ces
applications, avec l’attribution et la portée de chaque résultat.
